% ---
% title: Living Essays
% date: 2026-07
% category:
% nodes: tools 
% links: personal-website-design
% ---
\documentclass{article}
\usepackage[utf8]{inputenc}
\usepackage{amsmath, amssymb, amsthm}
\usepackage{geometry}
\usepackage{graphicx}
\usepackage{hyperref}

\geometry{a4paper, margin=1in}

\title{Living Essays}
\date{July 2026}
\author{Daksh Mehta}

\begin{document}
	
Future essays on this website will look slightly different. I am in the process of implementing a version history so viewers can see how an essay or a reflection has changed over time. I like the idea that writing and reflection is a continuous, not discrete, process. Essays wil now contain a slider with which readers can turn back time to view older versions of the current page. 

Most writing I've encountered whether on blog websites similar to this one, or academic papers seem like they landed as fully formed concepts. There are obviously many interesting drafts that get discarded and lost to the sands of time, deemed unworthy of final publication. I think these older drafts can be revealing and important too. Aside from lexical and grammatical changes, the narrative has natural perturbations as new evidence and ideas get folded into the discussion. 

In my opinion, there is also value in presenting a half-finished idea which can later be revisited. Not everything has to be polished with clear conclusions, most of the time actually it is unclear where to go and the development takes time. 

I've become more cognisant of this year, during my Master's. Throughout the year I tried to stay on top of writing small sections and paragraphs based on what I was doing at the time, in hopes that when I get to the write up at the end it will be easier to string my ideas together. Looking back through these old drafts reveals times I completely misunderstood something or an assumption I hadn't questioned. More commonly, the changes are subtle. An argument becomes a little more nuanced or a sentence becomes less certain. These changes are the honest record of how thinking develops and the finished product fails to capture this process. 

I’ve come to appreciate that knowledge is rarely static. Scientific understanding evolves through revision. Clinical practice changes as new evidence emerges. Good research is built on the willingness to revisit old conclusions in the light of better information.

This website has always been an experiment in capturing my thoughts and opinions at different points of time and how what I read and do influences them. This change attempts to become more authentic to the website's goal and also encourage me to write more without being worried about a final product. 
	
\end{document}
